\section{Conclusion}\label{sec:conclusion}

We believe our results open exciting research directions in both algorithm design and complexity theory, and we highlight a few here.

\paragraph{Fine-grained complexity.}
An important message from our paper is that the fine-grained reductions that FGC has built up are incredibly valuable, and that it is even more important to design more reductions to enable further algorithm design and improved conditional hardness.

{ Which hypotheses should take the place of $3$SUM and APSP?} As discussed in the introduction, one possibility is to restrict these hypotheses to combinatorial algorithms~\cite{VW10,VW18,AV14}, motivated by the search for more practical algorithms or practical hardness.
 Another possibility is the hypothesis that the balanced case of All-Edges Sparse Triangle, with $m$ edges, needs $m^{4/3-o(1)}$ time. Our algorithms do not speed up this balanced case, and since it is the target of several lower-bound reductions from $3$SUM and Exact Triangle~\cite{Pat10,KPP16,VX20a,CX24}, the hypothesis still gives many problems evidence of hardness. Is it possible that the lopsided version has a faster algorithm but the balanced version does not? 

More broadly, the complexity of the many problems that are only known to be 3SUM-hard or APSP-hard, rather than equivalent to $3$SUM or APSP, is now open (the gray boxes of Figure~\ref{fig:web}); we get no faster algorithms for them, since the reductions go the other way. For instance, can one decide in truly subquadratic time whether $n$ points in the plane contain three on a line~\cite{GO95}? As far as we know, the fastest known running time is still $O(n^2)$, without even any polylogarithmic improvements.  A similar question arises for the hypotheses that our results say nothing about, such as $k$-SUM for $k\ge4$, 3SUM-Indexing, OMv (without hints), and Orthogonal Vectors. Do ideas like ours help to speed up algorithms for any of them?

Our new technique gives a way to deal with sparsity, at least in lopsided instances. Much work has gone into algorithms for sparse matrix multiplication~\cite{YZ05,AP09,PS14,ABFK24,BGGW25,BFN26,Gra26}. However, our techniques do not seem to give new algorithms for sparse matrix multiplication. Is improving upon the known algorithms for this problem hard, or is there another technique that could help here? This question is related to the balanced version of All-Edges Sparse Triangle.

What are the true exponents of $3$SUM, APSP, and the many other problems of Figure~\ref{fig:web} that now have faster algorithms? Our exponents can certainly be improved, and one place to gain is the reductions. Each of our algorithms composes known reductions with the matrix theorem, and many reductions lose a constant fraction of the saving in the exponent. For instance, the reduction from Exact Triangle to Lopsided All-Edges Sparse Triangle keeps only half of the saving in the exponent, and the reductions from $3$SUM and APSP to Exact Triangle keep a half and a third (Section~\ref{sec:reduction}). These reductions, like most fine-grained reductions, were designed to prove hardness, where it only matters that they keep some polynomial saving. Now that they are being used as algorithms, how much they keep is more important. Which of them can be improved? Are there more efficient reductions from all these different problems to Lopsided All-Edges Sparse Triangle, or to thin matrix products themselves, that skip the intermediate problems?\footnote{\label{footnote:improvements}For one example, the exponents for APSP and $3$SUM presented here can already be improved slightly without too much work, by not going through the decision version of Exact Triangle. For APSP, the proof of Theorem~\ref{thm:red} also solves the all-edges version of Exact Triangle (scanning each pair only until its zero triangle is found), to which the $(\min,+)$-product reduces on the same $n$ vertices~\cite{VW10,VW18,VW13}, so that APSP keeps all of the saving of Exact Triangle rather than a third. For $3$SUM, the direct reduction of~\cite{KPP16} to offline Set Disjointness keeps all of the saving rather than a quarter, but it is randomized (its derandomization~\cite{FKP24} keeps far less of the saving). Substantially better improvements than these would be more exciting. We focused in this paper on simplifying the presentation rather than giving slight further improvements to the exponents.} 

Some of these losses cannot be removed by reductions alone. Sheffield, Vassilevska Williams, and Xi~\cite{SVX27} show that the loss of a third in the reduction from the all-edges version of a triangle problem to its detection version~\cite{VW10,VW18} is optimal for black-box reductions, those that work for every relation on the three edge weights at once, so a tighter reduction must use the structure of the particular problem. The loss can still be avoided on the algorithmic side, as in the footnote: our algorithm for Exact Triangle solves the all-edges version at no extra cost, as all known algorithms for triangle problems do~\cite{SVX27}. Our reductions are not black-box in this sense either: the reduction from Exact Triangle to Lopsided All-Edges Sparse Triangle hashes the weights modulo a prime, which uses the additive structure of the problem, so the barriers of~\cite{SVX27} say nothing against improving the half that it loses. The same paper also gives new black-box reductions among triangle, listing, and matrix problems, and such reductions are now of algorithmic interest as well.


Beyond the reductions, all of our algorithms use the same matrix theorem, and a different technique might give improvements such as better exponents, combinatorial or practical algorithms, or algorithms for All-Edges Sparse Triangle with a middle part larger than $n^{0.12}$, such as the $n^{1/2}$ that arises in an approach to girth approximation~\cite{RV12}.

\paragraph{Algebraic algorithm design.}
Any improvement to the matrix theorem carries over to all the other speedups we achieve here. We do not know the true complexity of computing a set $W$ of entries of a thin product $XY$, as a function of $D$ and $|W|$. In particular, to handle $D>N^{0.12}$, one could look for new base identities that compute longer inner products per multiplication than Sch\"onhage's, perhaps by computer search~\cite{FBH+22,KM23,Nov25}, or ask whether other known identities, such as border rank expressions for the Coppersmith--Winograd tensor~\cite{CW90},  which underlie the current bounds on $\omega$ and $\alpha$~\cite{VXXZ24,ADVXXZ25,DEK+26}, can be used in the same way. It is not immediately clear how to make use of other tensors, since our analysis uses specific properties of Sch\"onhage's identity, such as its sparsity and the way the leaves of its recursion are shared among the output entries (Section~\ref{sec:sparse}). If other tensors cannot be used in this way, it would be interesting to understand why. For instance, even if the rank or border rank of a tensor is known, it might be that we need to find different rank or border rank expressions to use, similar to work on improving leading constants of matrix multiplication algorithms (see, e.g.,~\cite{BB15,KS20,BS19,AY25,MS25}). 
